\section*{Chapitre \ref{ch:rc:structural-credit-models} --- Modèles de crédit structurels}

\begin{solution}{exo:rc:structural-credit-models:1}
Capitaux propres 18.89, dette $50-18.89 = 31.11$, un spread de rendement de 103 points de base
($-\ln(31.11/40)/5-0.04$), et une probabilité risque-neutre de défaut $N(-d_2) = 23.5~\%$.
\end{solution}

\begin{solution}{exo:rc:structural-credit-models:2}
Les actifs suivent une diffusion~: pour faire défaut dans les mois qui viennent, ils doivent passer de bien au-dessus de la dette
à un niveau inférieur à celle-ci, ce qu'un chemin continu, avec une volatilité de 10--20~\% par an, ne fait
presque jamais. La probabilité et le spread s'annulent plus vite que linéairement à mesure que l'échéance
se raccourcit.
\end{solution}

\begin{solution}{exo:rc:structural-credit-models:3}
$\mathrm{DD} = (\ln(40.69/40)+(0.08-\tfrac12\times0.156^2))/0.156 = 0.54$, une \omterm{def:rc:reduced-form-credit:pd}{probabilité de
défaut} de 29.3~\%, supérieure aux 15.4~\% à cinq ans~: sur un an, la dérive a eu peu de temps pour
soulever les actifs loin du nominal, et traiter tout le nominal comme exigible dans un an
surestime le risque. Les mesures pratiques de \omterm{def:rc:structural-credit-models:dd}{distance au défaut} retiennent un point de défaut
égal à la dette à court terme plus une partie de la dette à long terme, pour cette raison.
\end{solution}

\begin{solution}{exo:rc:structural-credit-models:4}
Avec $\nu = 0$, la formule se réduit au principe de réflexion~: $2N(\ln1.5/(0.3\sqrt2))-1 =
66.1~\%$.
\end{solution}

\begin{solution}{exo:rc:structural-credit-models:5}
La couverture est linéaire en notionnel~: le double de 1.74~millions de dollars, soit une vente à découvert
de 3.49~millions de dollars d'actions.
\end{solution}

\begin{solution}{exo:rc:structural-credit-models:6}
$\sigma_E = N(d_1)\sigma_VV/E$~: les capitaux propres sont une position à effet de levier sur les actifs, de sorte que
une variation en pourcentage des actifs se traduit par une variation en pourcentage plus grande des capitaux propres, dans le rapport
$N(d_1)V/E$ (3.2 pour l'entreprise du chapitre). Lorsque les capitaux propres montent, le levier recule et le rapport
baisse~: à volatilité des actifs donnée correspond une volatilité des capitaux propres qui diminue à mesure que le cours monte,
c'est l'effet de levier.
\end{solution}

\begin{solution}{exo:rc:structural-credit-models:7}
\texttt{implied\_barrier()} donne 78.7~\% du nominal~; le spread à dix ans du modèle avec cette
barrière vaut 380 points de base, moins que les 500 à cinq ans, car la courbe est en bosse.
\end{solution}

\begin{solution}{exo:rc:structural-credit-models:8}
On sait que le \omterm{def:rc:structural-credit-models:merton}{modèle de Merton} donne des spreads très inférieurs à ceux du marché pour la plupart des entreprises (pas de
défaut à court terme, pas de sauts, pas de liquidité ni de primes de risque au-delà de la diffusion), si bien que
un écart important est la norme, et non un signal. Un \omterm{def:rc:structural-credit-models:csa}{arbitrage de structure du capital} compare le marché à un
modèle calibré sur l'historique propre de l'entreprise (la barrière), et prend position sur les variations de l'écart,
couvert~; un arbitrage fondé sur le niveau exige encore une couverture et une opinion sur les raisons pour lesquelles l'écart
devrait se résorber.
\end{solution}

\begin{solution}{pb:rc:structural-credit-models:1}
\textbf{1.} 229 et 500 points de base~: le credit default swap (CDS) est 271 points de base au-dessus du
modèle, un signal pour vendre la protection.
\textbf{2.} $\partial s/\partial E = -0.0047$ par milliard de dollars~; vente à découvert de 1.74~millions
de dollars d'actions.
\textbf{3.} Le coupon à 500 points de base sur 10~millions de dollars (500\,000~dollars par an, avant le
coût de la vente à découvert d'actions et de ses dividendes, diminué du revenu de financement du produit de la vente).
\textbf{4.} Que le crédit de l'entreprise est bon marché par rapport à ses capitaux propres~: soit les spreads se resserrent,
ou les capitaux propres chutent.
\textbf{5.} Le défaut (un saut de 6~millions de dollars sur la protection contre une action
qui tombe presque à zéro), la convexité du modèle, la volatilité, les mouvements idiosyncratiques d'une
jambe (une OPA, une dégradation de notation), et la liquidité.
\textbf{6.} 161 points de base.
\textbf{7.} $500+(161-229) = 432$ points de base.
\textbf{8.} $-0.18\times1.74 \text{ million} = -\text{dollars}~313\,961$.
\textbf{9.} $+\text{dollars}~259\,002$ sur le CDS, un P\&L couvert de $-\text{dollars}~54\,958$
(la convexité de la \cref{fig:rc:structural-credit-models:equityspread}).
\textbf{10.} $-\text{dollars}~689\,903$ sur le CDS~; un total de $-\text{dollars}~1\,003\,863$.
\textbf{11.} $-0.0030$ par milliard de dollars~: la pente s'aplatit lorsque les capitaux propres montent, si bien que la couverture
doit être réduite.
\textbf{12.} La volatilité des actifs monte à 17.3~\% et le spread du modèle passe à 251 points
de base~: dans cette variante, le modèle prévoyait un spread plus large, et non plus serré, après le rallye.
\textbf{13.} Une OPA financée par la dette, ou un acheteur d'actions dont les intérêts diffèrent
de ceux des créanciers, fait monter à la fois la valeur des capitaux propres et le risque de crédit~: la valeur passe des
créanciers aux actionnaires, ce qu'un modèle à dette fixe et à valeur d'actif unique ne peut
produire.
\textbf{14.} 418 points de base~: le CDS aurait dû se resserrer de 82 points
de base pour compenser la perte de la jambe actions.
\textbf{15.} Rien dans la structure du capital ne couvre une dégradation tirée par la notation
elle-même~; une protection sur indice ou sur des noms voisins couvre une partie du risque, et une
position plus petite le limite.
\textbf{16.} Beaucoup détenaient des positions semblables dimensionnées par des modèles semblables, si bien que leurs dénouements
ont fait évoluer les cours en sens opposé, comme l'a décrit la revue de la BRI.
\textbf{17.} Par la perte en scénario de stress lorsque les deux jambes évoluent ensemble contre la position, et non par le
risque couvert du modèle~; avec des limites par nom et sur le risque d'événement.
\textbf{18.} Clôturer la position si l'événement a changé l'entreprise (la dette sera encore dégradée, l'investisseur
a l'intention de modifier la structure du capital~)~; conserver si l'écart reflète des ventes
forcées, et que l'entreprise est inchangée, avec de quoi financer l'attente.
\textbf{19.} Opération nommée~: \emph{l'\omterm{def:rc:structural-credit-models:csa}{arbitrage de structure du capital} de mai 2005}, où la jambe actions
perd 313\,961~dollars et le CDS 689\,903~dollars, soit au total 1\,003\,863~dollars sur
10~millions de dollars de protection sur deux jours.
\textbf{20.} Un lien entre les cours des capitaux propres et du crédit d'une même entreprise sert aux couvertures,
aux signaux et aux noms dépourvus de CDS liquides.
\end{solution}

\begin{solution}{iq:rc:structural-credit-models:1}
Les actionnaires ont une responsabilité limitée~: à l'échéance de la dette, ils la remboursent et gardent le reste,
ou renoncent et laissent les actifs aux créanciers. Leur gain vaut $\max(V_T-D,0)$, soit une option d'achat sur
les actifs, de prix d'exercice égal à la dette. La dette vaut $\min(V_T,D)$~: une obligation sans risque moins un put sur les
actifs.
\iqlookfor{responsabilité limitée, le call et la dette vue comme obligation moins put.}
\end{solution}

\begin{solution}{iq:rc:structural-credit-models:2}
Deux équations à deux inconnues~: la valeur des capitaux propres comme call sur les actifs, et la
volatilité des capitaux propres, égale à $N(d_1)\sigma_VV/E$~; on résout par Newton. Les autres méthodes itèrent sur la série des rendements
d'actifs déduite des valeurs quotidiennes des capitaux propres. Le point de défaut et l'horizon sont des choix, et
la volatilité des capitaux propres peut provenir des options.
\iqlookfor{les deux équations et les choix qui les sous-tendent.}
\end{solution}

\begin{solution}{iq:rc:structural-credit-models:3}
Les trajectoires continues des actifs ne peuvent pas atteindre rapidement le point de défaut, si bien que les \omterm{def:rc:reduced-form-credit:pd}{probabilités de défaut} à court terme
sont proches de zéro. Remèdes~: des sauts dans les actifs, une \omterm{def:rc:structural-credit-models:fp}{barrière de défaut} incertaine
(l'entreprise ne sait pas exactement où elle fait défaut), une information incomplète sur les
actifs, ou un modèle hybride avec une intensité de défaut.
\iqlookfor{la raison liée à la diffusion et un ou deux remèdes.}
\end{solution}

\begin{solution}{iq:rc:structural-credit-models:4}
Le nombre d'écarts-types d'actifs séparant la valeur espérée des actifs du point de défaut
à un horizon donné. Les praticiens le transposent en fréquences de défaut à l'aide d'une table empirique
construite à partir des défauts historiques, car la queue de la loi normale sous-estime les taux de défaut, et ils retiennent un
point de défaut situé entre la dette à court terme et la dette totale.
\iqlookfor{la formule et la transposition empirique.}
\end{solution}

\begin{solution}{iq:rc:structural-credit-models:5}
Vendre la protection et vendre à découvert les capitaux propres lorsque le swap de défaut est plus large que ne l'indique un \omterm{def:rc:structural-credit-models:structural}{modèle structurel}
calibré sur l'entreprise (ou l'inverse), le montant en capitaux propres étant fixé par
la sensibilité du spread aux capitaux propres, multipliée par le \omterm{def:rc:reduced-form-credit:cs01}{CS01} de la position. Elle perd sur les deux jambes si
les capitaux propres s'envolent (la vente à découvert perd) tandis que les spreads s'élargissent (la protection perd),
comme lors d'un rachat par emprunt (leveraged buyout) ou d'une offre publique d'achat, opérations qui font passer de la valeur des
créanciers vers les actionnaires, ou d'une dégradation qui ne frappe que la dette.
\iqlookfor{le ratio de couverture et le scénario où les deux jambes perdent.}
\end{solution}

\begin{solution}{iq:rc:structural-credit-models:6}
Les \omterm{def:rc:reduced-form-credit:reduced}{modèles à forme réduite} ajustent exactement la courbe de crédit et valorisent les dérivés de crédit
de façon cohérente~; ils ne disent rien des raisons pour lesquelles le défaut survient. Les \omterm{def:rc:structural-credit-models:structural}{modèles structurels} expliquent
le défaut à partir du bilan et relient le crédit aux capitaux propres, mais ajustent mal les spreads de marché.
Utiliser la forme réduite pour la valorisation et la gestion du risque des instruments de crédit~; le \omterm{def:rc:structural-credit-models:structural}{modèle structurel} pour les émetteurs
sans cotation de swaps de défaut, la valeur relative capitaux propres–crédit et le scoring de crédit.
\iqlookfor{ajustement contre explication, et l'usage de chacun.}
\end{solution}
